% ─────────────────────────────────────────────────────────────────────────────
% DISCUSSION (condensed — extended limitations in Supplementary Material)
% ─────────────────────────────────────────────────────────────────────────────

\section{Discussion}
\label{sec:discussion}

\subsection{What the Validation Does and Does Not Show}
\label{sec:discuss_validation}

The independent checks in \S\ref{sec:validation}--\ref{sec:iss_crosscheck}
(shielding curve shape, material ratio, the real MSL/RAD cruise trend, and
the real ISS multi-year trend) each probe an aspect of the model not
constrained by the single-point MSL/RAD calibration. Together they support
characterizing the pipeline as \emph{calibrated to one flight measurement
and consistent with independent benchmarks within their stated
uncertainties}. The
near-zero daily-resolution correlation, the SEP-event confound behind it,
and the held-out cross-validation of the RAD anchoring (\S\ref{sec:validation})
are reported for the same reason: an honest, non-circular account of what
this monthly-resolution model can and cannot claim is worth more than an
artificially strong number. That discipline paid off here: chasing the
daily correlation initially suggested a resolution gap in the solar-modulation
physics, but a genuine daily-resolution reconstruction (real neutron-monitor
data, independently calibrated) did not improve it, and the actual
explanation turned out to be a documented SEP event contaminating the
comparison for a few days.

\subsection{Implications for Mission Shielding Design}
\label{sec:discuss_shielding}

Turning from validation to what these numbers imply operationally, the
shielding curve (Fig.~\ref{fig:shielding_scan}) plateaus above
$\sim\!20\gcm$ aluminum as secondary-neutron growth offsets HZE
attenuation: going from $16$ to $40\gcm$ cuts GCR dose only
$\sim\!42\%$ while roughly doubling wall mass, consistent with prior
literature \citep{cucinotta2006, slaba2014} that bulk aluminum shielding
is mass-inefficient against chronic GCR. SEP shielding is the opposite
case: the October~2003 event drops from $\sim\!35\times$ the BFO limit
unshielded to below it by $\sim\!25\gcm$ (Fig.~\ref{fig:sep}). This
asymmetry supports a compact, heavily-shielded SEP storm shelter over
whole-vehicle HZE shielding as the more mass-efficient architecture
\citep{durante2011, cucinotta2010}.

\subsection{REID Interpretation: Biology Dominates}
\label{sec:discuss_reid}

The shielding trade-offs above are only as meaningful as the risk number
they are designed around, so it is worth scrutinizing that number directly.
Median REID ($6.5\%$ male, $8.3\%$ female at $16\gcm$) exceeds the
historical NASA $3\%$ limit \citep{nasa2014}, though NASA's updated
age/sex-dependent framework \citep{nasa2023} is not implemented here and
the comparison should be read as indicative. Because the
composition-constrained pipeline overestimates $\Qeff$ by $\sim\!27\%$
(\S\ref{sec:calib}), these values are conservative upper estimates;
correcting for that ratio would put male median REID at $\sim\!5.1\%$ —
still above the historical limit. The p5--p95 band spans a factor of
$\sim\!12\times$: this REID should be read as "the mission dose likely
carries meaningful cancer risk, with the exact value uncertain by an
order of magnitude."

Most importantly, this uncertainty is overwhelmingly biological. The four
radiobiological parameters contribute $\gtrsim\!85\%$ of REID's normalized
sensitivity, all five physics parameters combined $<15\%$ (Supplementary
Table~S4) — consistent with Cucinotta et al.\ \cite{cucinotta2013} and the
broader literature \citep{durante2011}. No direct human data constrain the
HZE LET range that dominates unshielded GCR dose equivalent; extrapolating
from mouse and atomic-bomb-survivor data to this regime is the single
largest source of uncertainty in this calculation, and no amount of
improved transport physics can resolve it. Strikingly, \code{phi\_scale}
drives $79\%$ of absorbed-dose variance but only $8\%$ of REID variance:
solar-modulation uncertainty strongly affects physical dose but has almost
no leverage on mortality probability once biology enters. The
value-of-information sweep (\S\ref{sec:vov}) makes this actionable rather
than qualitative.

\subsection{Scope of the LET/RBE Layer}

Separately from the REID interpretation above, it is worth being explicit
about what the pipeline's LET/RBE layer does and does not claim. The
per-species LET decomposition (Fig.~\ref{fig:let}) traces $\Qeff$'s
physical origin to specific ion contributions; the residual $27\%$
$\Qeff$ excess over the RAD measurement is attributed to CSDA's missing
fragmentation (\S\ref{sec:calib}). The repository also contains a proton
RBE self-consistency layer (\S\ref{sec:rbemethods}) reported only as an
implementation check (Table~\ref{tab:validation}); it is the foundation
for a planned, separate proton-therapy crossover study.

\subsection{Limitations}
\label{sec:discuss_limits}

The main structural limitation is CSDA transport without nuclear
fragmentation, which by my order-of-magnitude estimate overestimates
$\Qeff$ (and correspondingly REID) at $16\gcm$ by roughly
$15$--$25\%$ \citep{wilson1991, slaba2014} — well within the $12\times$
uncertainty band but not negligible relative to the median. Other
significant approximations include: a uniform-slab spacecraft geometry
(no ray-tracing); six modeled GCR species (unmodeled species $<5\%$ of
dose); the force-field modulation's monthly resolution, which limits
daily-scale claims and is the reason the ISS cross-check is restricted to
trend and does not extend to near-Earth or trapped-radiation
environments; the pre-2023 Cucinotta NASA risk framework rather than the
updated age/sex-dependent limits; and cancer-only risk (no CNS,
cataract, or cardiovascular endpoints). The full itemized discussion of
each limitation, with supporting estimates, is provided in Supplementary
\S{}S2.

\subsection{Future Work}
\label{sec:future_work}

Each limitation above points to a concrete next step, roughly in order of
expected impact on REID. \textbf{Nuclear fragmentation}: replacing the
Bradt-Peters fluence-removal model with tabulated fragmentation
cross-sections \citep{wilson1991} would let secondary light-ion production
be tracked explicitly, closing the $15$--$25\%$ $\Qeff$ overestimate
(\S\ref{sec:discuss_limits}) rather than only bounding it. \textbf{Risk
framework}: adopting NASA's 2023 age/sex-dependent career-limit framework
\citep{nasa2023} in place of the pre-2023 Cucinotta model would remove the
main caveat on the REID-vs-NASA-limit comparison in
\S\ref{sec:discuss_reid}. \textbf{Geometry}: ray-traced, non-uniform
spacecraft shielding (as in OLTARIS or Geant4) would replace the
uniform-slab approximation, most relevant for real vehicle interiors with
mass concentrated unevenly around the crew. \textbf{Non-cancer endpoints}:
CNS, cataract, and cardiovascular risk models are not yet implemented; CNS
effects in particular have rodent-study support for a cognitive-impairment
threshold below the cancer risk threshold modeled here. \textbf{Proton
RBE}: the dose-averaged LET/RBE self-consistency layer
(\S\ref{sec:rbemethods}) exists as the foundation for a planned, separate
proton-therapy crossover study.

\subsection{Comparison with Existing Tools and Reproducibility}
\label{sec:discuss_comparison}

Set against the alternatives, Table~\ref{tab:tool_comparison} compares
this pipeline against HZETRN, OLTARIS, Geant4, and TOPAS: it trades
high-fidelity transport physics for being, to my knowledge, the only
one that is simultaneously open/pip-installable and includes
organ-specific REID, an SEP library, and LHS uncertainty quantification.

\begin{table}[t]
  \centering
  \caption{Comparison of GCR dosimetry tools. \checkmark\ = present,
           $\triangle$ = partial, $\times$ = absent.}
  \label{tab:tool_comparison}
  \small
  \setlength{\tabcolsep}{4pt}
  \begin{tabular*}{\tblwidth}{@{}lccccc@{}}
    \toprule
    Feature & Ours & HZETRN & OLTARIS & Geant4 & TOPAS \\
    \midrule
    Open / pip-install   & \checkmark & $\times$    & $\times$    & \checkmark  & \checkmark \\
    Fragmentation        & $\times$   & \checkmark  & \checkmark  & \checkmark  & \checkmark \\
    Organ REID           & \checkmark & $\triangle$ & \checkmark  & $\times$    & $\times$   \\
    SEP library          & \checkmark & $\triangle$ & \checkmark  & $\times$    & $\times$   \\
    LHS uncertainty      & \checkmark & $\times$    & $\times$    & $\times$    & $\times$   \\
    Clinical RBE         & \checkmark & $\times$    & $\times$    & $\triangle$ & \checkmark \\
    CI test suite        & \checkmark & $\times$    & $\times$    & $\times$    & $\times$   \\
    \bottomrule
  \end{tabular*}
\end{table}

All results reproduce via \texttt{pip install -e .[dev]},
\texttt{validate\_pipeline.py}, \texttt{validate\_extended.py}, and
\texttt{pytest}; the full $N=500$ ensemble takes $\sim\!90$ minutes
single-core, with a reduced $N=20$ path used in CI. No proprietary
datasets are required.
